\subsection{Properties of the integral}

PROPOSITION 1. --- For every positive \(\mu\) -integrable numerical function \(f\) ,

\[
\int f d | \mu | = \int^ {*} f d | \mu | = \mathrm{N} _ {1} (f) \geqslant 0.\tag{2}
\]

For, \(\int f d|\mu|\) and \(\mathrm{N}_{1}(f)\) are continuous on \(L^{1}\) and are equal for every continuous function \(f \geqslant 0\) with compact support; on the other hand, every function \(f \geqslant 0\) in \(L^{1}\) is the limit (in the sense of convergence in mean) of a sequence of continuous functions \(\geqslant 0\) with compact support ( \(\S3\) , No.~5, Prop. 11); whence the proposition.

COROLLARY 1. --- For every integrable function \(\mathbf{f} \in \mathcal{L}_{\mathrm{F}}^{1}\) , \(|\mathbf{f}|\) is integrable and

\[
\int | \mathbf {f} | d | \mu | = \int^ {*} | \mathbf {f} | d | \mu | = \mathrm{N} _ {1} (\mathbf {f}).\tag{3}
\]

We shall make frequent use of Prop. 1 and its Cor. 1, on replacing \(\int^{*}fd|\mu|\) or \(\mathrm{N}_{1}(f)\) by \(\int fd|\mu|\) when dealing with an integrable function \(\geqslant 0\) . For example, for two integrable functions \(\mathbf{f},\mathbf{g}\) to be equivalent, it is necessary and sufficient that \(\int |\mathbf{f} - \mathbf{g}|d|\mu| = 0\) .

We recall that, for a function \(\mathbf{f}\) to belong to \(\mathcal{L}_{\mathrm{F}}^{p}\) , it is necessary and sufficient that the function \(|\mathbf{f}|^{p - 1}\cdot \mathbf{f}\) belong to \(\mathcal{L}_{\mathrm{F}}^{1}\) (§3, No.~8, Cor. 1 of Th. 7), that is, that it be integrable; this is the reason for the terminology ' \(p\) -th power integrable function'. Moreover:

COROLLARY 2. --- For every function \(\mathbf{f} \in \mathcal{L}_{\mathrm{F}}^{p}\) , the numerical function \(|\mathbf{f}|^{p}\) is integrable and

\[
\mathrm{N} _ {p} (\mathbf {f}) = \left(\int | \mathbf {f} | ^ {p} d | \mu |\right) ^ {1 / p}.\tag{4}
\]

This follows at once from the fact that \(|\mathbf{f}|\) belongs to \(\mathcal{L}^p\) (§3, No.~5, Prop. 11) and formula (2).

PROPOSITION 2. --- For every integrable function f,

\[
\left| \int \mathbf {f} d \mu \right| \leqslant \int | \mathbf {f} | d | \mu |.\tag{5}
\]

This follows at once from the inequality (1) by passage to the limit, on taking into account (3) and the continuity of \(\mathrm{N}_{1}(\mathbf{f})\) on \(L_{F}^{1}\) .

THEOREM 1. --- Let F and G be two Banach spaces, u a continuous linear mapping of F into G. For every integrable function f with values in F, \(u \circ f\) is integrable and

\[
\int \mathbf {u} (\mathbf {f} (x)) d \mu (x) = \mathbf {u} \left(\int \mathbf {f} (x) d \mu (x)\right).\tag{6}
\]

We already know that \(\mathbf{u} \circ \mathbf{f}\) is integrable (§3, No.~5, Th. 4); the relation (6), being valid for every \(\mathbf{f} \in \mathcal{K}_{\mathrm{F}}\) , extends to every integrable function \(\mathbf{f}\) by the principle of extension of identities: for, \(\mathbf{f} \mapsto \mathbf{u} \circ \mathbf{f}\) is continuous for the topology of convergence in mean, as follows from the inequality \(\mathrm{N}_1(\mathbf{u} \circ \mathbf{f}) \leqslant \| \mathbf{u} \| \cdot \mathrm{N}_1(\mathbf{f})\) .

COROLLARY 1. --- Let \(\mathbf{a}'\) be a continuous linear form on \(\mathbf{F}\) . If \(\mathbf{f}\) is an integrable function with values in \(\mathbf{F}\) , then the numerical function \(\langle \mathbf{f}, \mathbf{a}' \rangle\) is integrable and

\[
\int \left\langle \mathbf {f} (x), \mathbf {a} ^ {\prime} \right\rangle d \mu (x) = \left\langle \int \mathbf {f} (x) d \mu (x), \mathbf {a} ^ {\prime} \right\rangle .\tag{7}
\]

We shall see in Ch. VI, §1, Exers. 7, 11 and 12 that there can exist functions \(\mathbf{f}\) , with values in an infinite-dimensional Banach space \(\mathbf{F}\) , such that \(\langle \mathbf{f}, \mathbf{a}' \rangle\) is integrable for every continuous linear form \(\mathbf{a}'\) on \(\mathbf{F}\) , without \(\mathbf{f}\) being integrable.

COROLLARY 2. --- If the \(\mathbf{a}_k\) ( \(1 \leqslant k \leqslant n\) ) are vectors in F and the \(f_k\) ( \(1 \leqslant k \leqslant n\) ) are integrable numerical functions, then the function \(\mathbf{f} = \sum_{k=1}^{n} \mathbf{a}_k f_k\) is integrable and

\[
\int \left(\sum_ {k = 1} ^ {n} \mathbf {a} _ {k} f _ {k}\right) d \mu = \sum_ {k = 1} ^ {n} \mathbf {a} _ {k} \int f _ {k} d \mu .\tag{8}
\]

